\projectentry
  {Tungsten}
  {Open-source scientific-computing plugin for Neovim}
  {Personal software project}
  {2025 -- Present}
  {\portfoliotag{Lua} \portfoliotag{Neovim} \portfoliotag{\LaTeX} \portfoliotag{Wolfram Engine} \portfoliotag{Python} \portfoliotag{Scientific Computing}}
  {\tungstenfigure}
  {Tungsten brings a computational-notebook workflow into ordinary \LaTeX\ editing inside Neovim. Instead of switching between editor, CAS, and scratch files, mathematical expressions can be interpreted and evaluated in place, with results rendered directly in the buffer.

  \begin{projectbullets}
    \item Designed and implemented a general-purpose LaTeX math parser covering algebra and arithmetic, calculus, linear algebra, differential equations, plotting, numerics, and units.
    \item Broke a large, under-defined parsing problem into manageable components, gradually building support for real mathematical notation such as systems of equations, matrices, vectors, nested expressions, Greek symbols, assumptions, units, and plotting commands.
    \item Integrated the parser into a Neovim workflow with inline results, user-facing commands, and backend support for Wolfram Engine and Python.
    \item Maintained the project through releases, documentation, dependency management, issue-driven improvements, and feedback from external users; currently has 20 GitHub stars, 36 user-facing commands, and 180 unique cloners in the last 14 days.
  \end{projectbullets}}

\projectseparator

\projectentry
  {Guitar Pedal Circuit Design}
  {Analog audio circuit design with gain, clipping, and tone-shaping stages}
  {Personal electronics project}
  {2025 -- Present}
  {\portfoliotag{Analog electronics} \portfoliotag{KiCad} \portfoliotag{ngspice} \portfoliotag{Tone shaping} \portfoliotag{Oscilloscope testing}}
  {\pedalfigure}
  {This project explores the design of a highly customizable distortion pedal from the ground up. Rather than modifying an existing circuit, the design is built as a complete analog signal chain with power conditioning, buffering, gain staging, clipping, clean blend, EQ, tone shaping, and output buffering.

  \begin{projectbullets}
    \item Designed an almost 100-component analog circuit in KiCad, including power supply section, input buffer, two gain stages, voice filtering, selectable clipping, clean blend, tilt control, volume control, and output buffer.
    \item Built the clipping section around six selectable diode options, with switching that allows either symmetrical clipping of both waveform halves or clipping of only one half.
    \item Used ngspice to inspect operating points, transient response, and frequency response before moving the circuit to breadboard testing.
    \item Breadboarded and tested sections of the circuit using measurements and oscilloscope inspection, with perfboard layout planning and enclosure CAD as the next steps.
  \end{projectbullets}}

\projectseparator

\clearpage

\projectentry
  {DC Motor Flap Control System}
  {Arduino-based PID position control for an aircraft wing flap model}
  {Control systems project}
  {BSc project, Ongoing}
  {\portfoliotag{PID control} \portfoliotag{Arduino} \portfoliotag{MATLAB} \portfoliotag{Simulink} \portfoliotag{System identification}}
  {\flapfigure}
  {This project develops a closed-loop control system for a 3D-printed aircraft wing flap model driven by a geared DC motor. The angular position of the motor shaft corresponds to the flap angle, and the objective is to achieve reliable position control with strong disturbance rejection, fast settling, and limited overshoot.

  \begin{projectbullets}
    \item Building a physical test setup using a 6–12 V DC motor driven from a 9 V supply, an integrated encoder wired to an Arduino Uno R3 WiFi, and a 3D-printed airfoil/flap model.
    \item Identifying an experimental transfer function from measured motor and encoder data to support controller design and validation.
    \item Designing and tuning a custom PID controller using root locus, Bode plots, gain/phase margins, and time-domain response analysis.
    \item Combining Arduino/C++, MATLAB, Simulink, Control System Toolbox, CAD, and 3D printing to connect modelling, implementation, and physical testing.
  \end{projectbullets}}

\projectseparator

\projectentry
  {Audio-Reactive TV Backlight}
  {Analog LED-control circuit driven by a TV audio signal}
  {Electromagnetism capstone project}
  {BSc project, 2026}
  {\portfoliotag{Analog electronics} \portfoliotag{Circuitry} \portfoliotag{Signal conditioning} \portfoliotag{Falstad} \portfoliotag{KiCad}}
  {\leddriverfigure}
  {This project involved designing, simulating, and building an analog circuit that converts a TV audio signal into a smooth LED-strip brightness response. The goal was not just to detect audio level, but to create a visually pleasing light response that followed the general energy and mood of the soundscape.

  \begin{projectbullets}
    \item Designed the full analog signal chain: TV audio output through SCART/RCA adapter, mono summing, AC coupling, non-inverting amplification and rectification, low-pass smoothing, current sink, and LED-strip output.
    \item Used Falstad simulations to test and tune the circuit before soldering the final version on perfboard.
    \item Calculated and adjusted gain and filtering behaviour to convert the audio into an envelope suitable for LED brightness control.
    \item Tuned the response through trial and error, balancing smoothness and responsiveness to avoid excessive flicker while still matching the perceived vibe of the audio.
  \end{projectbullets}}

\projectseparator

\projectentry
  {Game of Throws Launch Mechanism}
  {Spring-powered mechanical launcher for squash-ball competition}
  {Design \& Construction course project}
  {BSc project, 2024}
  {\portfoliotag{Mechanism design} \portfoliotag{SolidWorks} \portfoliotag{3D printing} \portfoliotag{Laser cutting} \portfoliotag{Testing}}
  {\catapultfigure}
  {This five-person course project involved designing and building a spring-powered mechanism for launching squash balls toward a target plate with scoring holes. The mechanism had to fit within strict packaging, material, loading-time, and adjustability constraints while delivering enough ball speed to be competitive.

  \begin{projectbullets}
    \item Worked within constraints including a restricted packaging volume, limited 3D-print volume, restricted materials, spring-powered launch energy, horizontal and vertical adjustment, and a maximum 4-second loading requirement.
    \item Contributed to the documentation and engineering reasoning behind the design, including requirements, concept evaluation, prototype testing, and design iteration.
    \item Helped communicate how testing revealed excessive friction and mechanical losses as key limitations in the first concepts.
    \item Documented the resulting pulley-based redesign, which aimed to extract more useful launch force from the available spring energy while preserving buildability and repeatable operation.
  \end{projectbullets}}
